% ============================================================
% CHƯƠNG 3: THIẾT KẾ KIẾN TRÚC HỆ THỐNG
% MỤC ĐÍCH: Đáp ứng Câu 4 (giải pháp + biện luận — 30 điểm, mục cao nhất)
% PHÂN BỔ: ~12–14 trang (chương quan trọng nhất)
% ============================================================
\chapter{Thiết kế kiến trúc hệ thống}
\label{chap:kien-truc}

Chương này trình bày cách ELDAS được thiết kế để giải bài toán đã đặt ra
ở Chương~\ref{chap:gioi-thieu}. Hệ thống được chia thành \textbf{bốn
lớp}; mỗi lớp đảm nhận một vai trò riêng và giao tiếp với nhau qua giao
diện được định nghĩa rõ ràng. Phần mở đầu giới thiệu tổng quan kiến
trúc cùng lý do lựa chọn từng công nghệ. Các phần kế tiếp lần lượt mô
tả chi tiết bốn lớp: mô phỏng (CloudSim Plus), cầu nối ngôn ngữ (Py4J),
tác tử học (MO-Gymnasium) và lớp đóng gói hạ tầng (Docker Compose).

% ============================================================
\section{Tổng quan kiến trúc bốn lớp}
\label{sec:tong-quan-kien-truc}

\subsection{Vì sao chia thành bốn lớp?}
ELDAS phải kết hợp ba mảng rất khác nhau: \textbf{mô phỏng cluster}
(CPU, RAM, GPU, năng lượng), \textbf{thuật toán học} (PPO, mạng nơ-ron)
và \textbf{đóng gói triển khai} (Docker, volume, biến môi trường). Nếu
dồn tất cả vào một chương trình duy nhất, ngay lập tức sẽ vấp phải ba
khó khăn:
\begin{itemize}
  \item Engine mô phỏng năng lượng datacenter tốt nhất hiện nay
        (CloudSim Plus) chỉ có bản Java, trong khi các thư viện RL hiện
        đại (PyTorch, Stable-Baselines3, MO-Gymnasium) đều dùng
        Python --- hai ngôn ngữ không thể chạy trong cùng một process.
  \item Mỗi lần đổi siêu tham số huấn luyện không nên kéo theo việc
        khởi động lại cả engine mô phỏng, vì khâu này rất tốn thời gian.
  \item Người khác cần dựng lại được toàn bộ hệ thống chỉ bằng một câu
        lệnh, thay vì phải cài thủ công từng phần.
\end{itemize}

Lời giải là chia hệ thống thành bốn lớp
(Hình~\ref{fig:arch-overview}). Java và Python chạy trong hai container
Docker riêng, trao đổi với nhau qua một cầu nối trong cùng mạng nội bộ.

\begin{figure}[H]
  \centering
  \begin{tikzpicture}[
      font=\footnotesize,
      node distance=4mm and 6mm,
      box/.style    ={draw, rounded corners=2pt, minimum height=10mm,
                      minimum width=30mm, align=center, fill=white,
                      inner sep=2pt},
      l1/.style     ={box, fill=gray!10},
      l2/.style     ={box, fill=blue!8},
      l3/.style     ={box, fill=green!8, minimum width=58mm},
      l4/.style     ={box, fill=orange!10},
      arrow/.style  ={-{Latex[length=2mm]}, thick},
      bidir/.style  ={{Latex[length=2mm]}-{Latex[length=2mm]}, thick}
    ]
    % ── L4 (top): Python MORL ──
    \node[l4] (state)   {\texttt{state\_}\\\texttt{builder}};
    \node[l4, right=of state] (rew) {\texttt{reward}};
    \node[l4, right=of rew]   (env) {\texttt{environment}\\(MO-Gym)};
    \node[l4, right=of env]   (trk) {\texttt{tracker}\\(W\&B)};

    % ── L2: Java sim (đặt trước L3 để căn giữa L3 theo L2/L4) ──
    % Tăng khoảng cách từ 22mm lên 50mm để tạo không gian rộng rãi cho Lớp 3 ở giữa
    \node[l2, below=50mm of state] (sim)  {\texttt{Simulation}\\\texttt{Manager}};
    \node[l2, right=of sim]        (dcf)  {\texttt{Datacenter}\\\texttt{Factory}};
    \node[l2, right=of dcf]        (trc)  {\texttt{AlibabaTrace}\\\texttt{Reader}};
    \node[l2, right=of trc]        (pol)  {\texttt{VmAlloc}\\\texttt{Policy}};

    % ── L3: bridge — đặt giữa L4 và L2, căn theo trung tâm cột rew/env ──
    % Đẩy xuống -25mm để nằm chính giữa khoảng không gian 50mm
    \node[l3] (gw) at ($(rew.south)!0.5!(env.south) + (0,-25mm)$)
              {\texttt{Py4J Gateway} (TCP :25333)};

    % ── L1: Docker — căn giữa theo cụm L2 ──
    % Tăng khoảng cách từ 15mm lên 25mm để tách biệt khỏi Lớp 2
    \node[l1, minimum width=140mm] (dc)
          at ($(sim.south)!0.5!(pol.south) + (0,-25mm)$) {%
      \texttt{docker compose} \;\;|\;\; bridge \texttt{sim-net} \;\;|\;\;
      volumes: \texttt{trace-data}, \texttt{results}, \texttt{model-store}};

    % ── khung đứt nét cho 4 lớp (vẽ sau khi tất cả node đã đặt) ──
    % Tăng inner sep để tiêu đề không đè lên các node bên trong
    \node[fit=(state)(trk),
          draw, dashed, inner sep=7mm,
          label={[anchor=north west,xshift=2mm,yshift=-1mm,font=\small\bfseries]
                  north west:Lớp 4 — MO-Gymnasium Agent (Python 3.11)}] (L4) {};
    \node[fit=(gw),
          draw, dashed, inner sep=6mm,
          label={[anchor=north west,xshift=2mm,yshift=-1mm,font=\small\bfseries]
                  north west:Lớp 3 — Cầu nối Java\,$\leftrightarrow$\,Python}] (L3) {};
    \node[fit=(sim)(pol),
          draw, dashed, inner sep=7mm,
          label={[anchor=north west,xshift=2mm,yshift=-1mm,font=\small\bfseries]
                  north west:Lớp 2 — CloudSim Plus 8.5.7 (Java 21)}] (L2) {};
    \node[fit=(dc),
          draw, dashed, inner sep=6mm,
          label={[anchor=north west,xshift=2mm,yshift=-1mm,font=\small\bfseries]
                  north west:Lớp 1 — Docker Compose v2 (host OS)}] (L1) {};

    % ── arrows: tất cả đều dùng anchor cạnh hộp, không xuyên qua hộp khác ──
    % env (L4) ↔ gw (L3): thẳng đứng từ cạnh dưới env xuống cạnh trên gw
    \draw[bidir] (env.south) -- (env.south |- gw.north)
                  node[midway, right=1mm, font=\scriptsize]
                  {\texttt{step()/reset()}};
    % gw (L3) ↔ sim (L2): từ cạnh trái-dưới gw xuống cạnh trên-phải sim (đường chéo)
    \draw[bidir] (gw.south west) -- (sim.north east)
                  node[midway, above, sloped, font=\scriptsize]{Py4J RPC};
    % trace CSV: từ volume L1 (cạnh trên) lên cạnh dưới trc (L2)
    \draw[arrow] (dc.north -| trc.south) -- (trc.south)
                  node[midway, right=1mm, font=\scriptsize]{trace CSV};
\end{tikzpicture}
  \caption{Kiến trúc bốn lớp của ELDAS. Mỗi đường đứt nét là một biên cô lập
  về ngôn ngữ, container, hoặc hạ tầng. Mũi tên đậm đánh dấu kênh giao tiếp
  hai chiều của vòng lặp RL.}
  \label{fig:arch-overview}
\end{figure}

\subsection{Vai trò từng lớp}
\textbf{Lớp 1 --- Docker Compose.} Đảm bảo toàn bộ hệ thống chạy đồng
nhất trên mọi máy. Hai container (\texttt{cloudsim-java} và
\texttt{rl-agent}) dùng chung một mạng nội bộ cùng ba volume chứa dữ
liệu trace, kết quả và checkpoint.

\textbf{Lớp 2 --- CloudSim Plus.} Đóng vai trò ``môi trường'' theo
nghĩa của RL: chứa cluster ảo (host, CPU, RAM, GPU), đọc trace
Alibaba, nhận hành động từ agent và trả về trạng thái mới kèm phần
thưởng. Toàn bộ mô hình năng lượng và GPU đều nằm tại lớp này.

\textbf{Lớp 3 --- Py4J Gateway.} Đảm nhận vai trò cầu nối giữa Java và
Python. Lớp này chỉ phơi ra một vài phương thức cơ bản
(\texttt{reset}, \texttt{step}, \texttt{getActionMask},
\texttt{exportMetrics}). Phía Python không cần biết bất kỳ chi tiết
nào về CloudSim, nên sau này hoàn toàn có thể thay engine bằng KWOK
hoặc một cluster Kubernetes thật mà không phải đụng đến code Python.

\textbf{Lớp 4 --- MO-Gymnasium Agent.} Đóng vai trò ``tác tử'' RL:
nhận quan sát từ Java, dùng MaskablePPO để chọn hành động và ghi log
lên Weights \& Biases. Toàn bộ hệ sinh thái Python (PyTorch,
Stable-Baselines3, MO-Gymnasium) được giữ trọn trong lớp này.

\subsection{Một bước trong vòng lặp RL diễn ra như thế nào?}
Hình~\ref{fig:rl-sequence} minh hoạ quá trình của một lần
\texttt{step()}: Python gửi hành động sang Java, Java cập nhật mô
phỏng rồi trả về quan sát mới kèm phần thưởng. Toàn bộ một vòng
\emph{round-trip} chỉ tốn khoảng $0{,}1$--$0{,}3$\,ms khi hai
container cùng chạy trên một máy.

\begin{figure}[H]
  \centering
  \begin{tikzpicture}[
      font=\small,
      lifeline/.style={draw, very thick, dashed, gray!60},
      msg/.style={-{Latex[length=2mm]}, thick},
      block/.style={draw, fill=blue!8, minimum width=22mm, minimum height=6mm,
                    rounded corners=1pt}
    ]
    % three actors
    \node[block] (py) at (0,0)  {Agent Python};
    \node[block] (gw) at (5,0)  {Py4J Gateway};
    \node[block] (sm) at (10,0) {SimulationManager};

    \foreach \x in {0,5,10} \draw[lifeline] (\x,-0.4) -- (\x,-7.5);

    % messages
    \draw[msg] (0,-1)   -- node[above,font=\scriptsize]{\texttt{step(action)}}     (5,-1);
    \draw[msg] (5,-1.6) -- node[above,font=\scriptsize]{\texttt{actionQueue.put}}  (10,-1.6);

    \node[draw, fill=orange!10, minimum width=30mm, anchor=west, font=\scriptsize]
      at (8.5,-2.5) {allocateTask + advanceEnergy};
    \node[draw, fill=orange!10, minimum width=30mm, anchor=west, font=\scriptsize]
      at (8.5,-3.3) {computeReward + recordSnapshot};
    \node[draw, fill=orange!10, minimum width=30mm, anchor=west, font=\scriptsize]
      at (8.5,-4.1) {buildObservation};

    \draw[msg] (10,-5)   -- node[above,font=\scriptsize]{\texttt{stepResultQueue.put}} (5,-5);
    \draw[msg] (5,-5.6)  -- node[above,font=\scriptsize]{\texttt{StepResult} (obs, r, done)} (0,-5.6);

    \node[draw, fill=green!10, anchor=west, font=\scriptsize]
      at (-1.4,-6.5) {numpy.float32, action\_masks};

    \draw[msg] (0,-7.2)  -- node[above,font=\scriptsize]{next \texttt{action}}  (5,-7.2);
  \end{tikzpicture}
  \caption{Sequence diagram cho một vòng lặp RL. Hai
  \texttt{SynchronousQueue} ở phía Java đảm bảo simulation thread chỉ tiến
  bước khi nhận được hành động hợp lệ.}
  \label{fig:rl-sequence}
\end{figure}

% ============================================================
\section{Lựa chọn giải pháp và biện luận}
\label{sec:lua-chon-giai-phap}

Bốn quyết định thiết kế lớn nhất của ELDAS đều được trình bày theo cùng
một khuôn: nêu các phương án, so sánh ưu nhược điểm trên các tiêu chí
thực tế, rồi chốt lại phương án kèm phần đánh đổi cần chấp nhận.
Bảng~\ref{tab:decision-matrix} ở cuối phần này tổng hợp lại cả bốn quyết
định.

\subsection{Lựa chọn công cụ mô phỏng: KWOK vs.\ CloudSim Plus}
\label{ssec:choose-sim}

Hai ứng viên cho việc mô phỏng cluster có khả năng tạm dừng phục vụ RL
là CloudSim Plus và KWOK (\textit{Kubernetes WithOut Kubelet}). Có ba
điểm quyết định lựa chọn:

\begin{enumerate}
  \item \textbf{Tốc độ.} CloudSim Plus dùng cơ chế
        \emph{discrete-event simulation}, có thể tua nhanh qua các
        khoảng không có sự kiện; một episode $5\,000$ tác vụ chạy
        xong trong vài giây. KWOK chạy ở thời gian thực, cùng
        episode đó tốn tới nhiều phút --- chậm hơn hàng trăm lần và
        không phù hợp cho việc huấn luyện hàng triệu bước.
  \item \textbf{Mô hình năng lượng.} CloudSim Plus có sẵn mô hình
        công suất CPU tuyến tính theo SPECpower; trong khi KWOK
        không có khái niệm năng lượng và phải tự viết plugin để bổ
        sung.
  \item \textbf{Tạm dừng / tiếp tục.} Vòng lặp RL cần có khả năng
        dừng lại tại mỗi sự kiện ``tác vụ mới đến'' để chờ agent ra
        quyết định. CloudSim Plus hỗ trợ pause/resume sẵn có; KWOK
        thì không.
\end{enumerate}

\textbf{Quyết định:} chọn \textbf{CloudSim Plus 8.5.7}.
\textbf{Đánh đổi cần chấp nhận:} mô hình mạng còn đơn giản và không có
\texttt{kubectl} thật, nên về sau, nếu muốn kiểm chứng trên cluster
thực, vẫn cần thêm bước port chính sách. Tuy vậy, do Lớp 3 đã được
giữ rất gọn, việc chuyển đổi chỉ cần thay phần backend Java; phần code
Python không phải sửa.

\subsection{Lựa chọn cầu nối ngôn ngữ: Py4J vs.\ gRPC vs.\ REST}
\label{ssec:choose-bridge}

Bảng~\ref{tab:bridge-comparison} ở §2.5 đã so sánh chi tiết. Ở đây
nhắc lại ba lý do chính khiến Py4J phù hợp nhất với ELDAS:
\begin{itemize}
  \item \textbf{Gọn nhẹ.} Mỗi bước RL chỉ trao đổi vài chục số
        thực --- nhỏ hơn nhiều so với phần overhead header và bước
        parse JSON mà REST hay gRPC kéo theo.
  \item \textbf{Tần suất cao.} Phase 2 sẽ gọi \texttt{step()} hàng
        triệu lần. Với REST (khoảng $5$\,ms mỗi lần), riêng phần
        giao tiếp đã tiêu tốn vài giờ; trong khi Py4J (khoảng
        $0{,}1$\,ms) chỉ cần khoảng vài chục phút.
  \item \textbf{Truy cập trực tiếp đối tượng Java.} Py4J cho phép
        Python gọi thẳng phương thức của đối tượng Java, không cần
        định nghĩa \texttt{.proto} hay OpenAPI schema --- giảm
        đáng kể lượng boilerplate.
\end{itemize}

\textbf{Quyết định:} chọn \textbf{Py4J 0.10.9.7}.
\textbf{Đánh đổi cần chấp nhận:} hai container phải nằm trong cùng
mạng Docker, không thể tách ra hai máy khác nhau. ELDAS chấp nhận
giới hạn này vì phạm vi đồ án chỉ chạy trong phòng thí nghiệm.

\subsection{Lựa chọn thuật toán RL: PPO vs.\ DQN vs.\ A2C}
\label{ssec:choose-rl}

Mục~\ref{sec:morl} đã giới thiệu PPO. Ở đây, ba ứng viên được so sánh
trên ba tính chất bắt buộc của bài toán:
\begin{enumerate}
  \item \textbf{Hỗ trợ action masking.} Ở mỗi bước, không phải host
        nào cũng còn đủ tài nguyên cho task hiện tại, nên agent
        cần bị cấm cứng khỏi các host không khả thi. PPO có biến
        thể \texttt{MaskablePPO} trong \texttt{sb3-contrib} hỗ trợ
        sẵn; DQN phải tự sửa cả bước chọn action lẫn bước Bellman;
        còn A2C chưa có triển khai mask sẵn dùng.
  \item \textbf{Ổn định khi reward thưa và lệch biên độ.} Reward
        năng lượng luôn mang giá trị âm cỡ vài trăm, còn reward
        SLA gần như bằng $0$ ở hầu hết các bước. Trước phổ reward
        kiểu này, DQN dễ bị overestimate giá trị $Q$; trong khi
        PPO có hai cơ chế bù trừ là chuẩn hoá advantage và cắt tỉ
        số chính sách.
  \item \textbf{Dễ song song hoá.} Phase 2 sẽ huấn luyện nhiều
        agent cùng lúc, mỗi agent ứng với một trọng số khác nhau.
        PPO on-policy chạy nhiều môi trường song song một cách tự
        nhiên; trong khi DQN với replay buffer đòi hỏi kiến trúc
        phức tạp hơn nhiều.
\end{enumerate}

\textbf{Quyết định:} \textbf{PPO + MLP}, dùng \textbf{MaskablePPO} từ
\texttt{sb3-contrib}. \textbf{Hướng mở rộng:} có thể thay MLP bằng GNN
ở giai đoạn sau --- chỉ cần đổi backbone, không phải động đến bốn lớp
kiến trúc.

\subsection{Lựa chọn chiến lược MORL: Linear Scalarization vs.\ NSGA-II}
\label{ssec:choose-morl}

\textbf{Linear scalarization} chuyển bài toán đa mục tiêu thành nhiều
bài toán đơn mục tiêu: với mỗi trọng số $w$, agent học một chính sách
tối ưu tại điểm cân bằng tương ứng; quét nhiều trọng số rồi nối các
điểm thu được sẽ thành Pareto front. \textbf{NSGA-II}~\cite{deb2002fast}
ngược lại là thuật toán tiến hoá, sinh ra cả Pareto front trong một
lần chạy. Tuy vậy, xét trong bối cảnh của ELDAS:
\begin{itemize}
  \item NSGA-II phải đánh giá hàng nghìn cá thể ở mỗi thế hệ, và
        mỗi cá thể tương ứng với một episode mô phỏng hoàn chỉnh.
        Tổng chi phí vượt xa ngân sách thời gian của Phase 2.
  \item Linear scalarization chỉ cần huấn luyện 9 chính sách độc
        lập (ứng với 9 giá trị $w$) và dễ chạy song song trên
        nhiều CPU. Khi triển khai thực tế, mỗi quyết định chỉ tốn
        một lượt forward pass của mạng --- đủ nhanh cho online
        scheduling.
  \item Đúng là linear scalarization không khôi phục được những
        đoạn lõm của Pareto front, nhưng hạn chế này đã được nêu
        rõ ở §\ref{sec:morl} và có thể chấp nhận trong phạm vi đồ
        án.
\end{itemize}

\textbf{Quyết định:} dùng \textbf{linear scalarization} với weight
sweep $w \in \{0{,}1;\, 0{,}2;\, \dots;\, 0{,}9\}$. NSGA-II và các
thuật toán MORL chuyên dụng (Envelope-Q, Pareto-DQN) được giữ ở phần
hướng phát triển tương lai.

\subsection{Bảng tổng hợp ma trận quyết định}
\begin{table}[H]
  \centering
  \caption{Ma trận quyết định bốn lựa chọn kiến trúc của ELDAS. Ký hiệu:
  \checkmark\,= đáp ứng tốt, $\sim$\,= chấp nhận được, $\times$\,= không
  đáp ứng.}
  \label{tab:decision-matrix}
  \footnotesize
  \begin{tabular}{lccccccl}
    \toprule
    \textbf{Quyết định} & \textbf{Phương án} &
      \textbf{Latency} & \textbf{Năng lượng} & \textbf{Pause/Resume} &
      \textbf{Đa mục tiêu} & \textbf{Reproducible} & \textbf{Chốt} \\
    \midrule
    \multirow{2}{*}{Engine}   & CloudSim Plus     & \checkmark & \checkmark & \checkmark & \checkmark & \checkmark & $\bigstar$ \\
                              & KWOK              & $\sim$     & $\times$   & $\times$   & \checkmark & \checkmark &            \\ \midrule
    \multirow{3}{*}{Bridge}   & Py4J              & \checkmark & ---        & ---        & ---        & \checkmark & $\bigstar$ \\
                              & gRPC              & $\sim$     & ---        & ---        & ---        & \checkmark &            \\
                              & REST              & $\times$   & ---        & ---        & ---        & \checkmark &            \\ \midrule
    \multirow{3}{*}{RL Algo}  & PPO + MaskablePPO & ---        & ---        & ---        & \checkmark & \checkmark & $\bigstar$ \\
                              & DQN               & ---        & ---        & ---        & $\sim$     & \checkmark &            \\
                              & A2C               & ---        & ---        & ---        & $\times$   & \checkmark &            \\ \midrule
    \multirow{2}{*}{MORL}     & Linear scalar.    & ---        & ---        & ---        & $\sim$     & \checkmark & $\bigstar$ \\
                              & NSGA-II           & $\times$   & ---        & ---        & \checkmark & \checkmark &            \\
    \bottomrule
  \end{tabular}
\end{table}

% ============================================================
\section{Thiết kế Lớp 2 — CloudSim Plus}
\label{sec:layer-cloudsim}

Lớp 2 là phần lớn nhất của hệ thống, gồm khoảng $1\,800$ dòng Java
chia thành chín lớp (Hình~\ref{fig:class-diagram}). Mỗi lớp đảm nhận
đúng một công việc cụ thể, qua đó việc bảo trì và mở rộng đều dễ dàng
hơn.

\begin{figure}[H]
  \centering
  \begin{tikzpicture}[
      font=\scriptsize,
      node distance=6mm and 5mm,
      cls/.style={draw, rectangle split, rectangle split parts=2,
                  rectangle split part align=base, fill=blue!5,
                  text width=30mm, align=left, rounded corners=1pt,
                  inner sep=2pt},
      arrow/.style={-{Latex[length=1.8mm]}, thick, gray!70}
    ]
      % ── Hàng 1: Gateway ──
      \node[cls, fill=green!8] (gw) {\textbf{GatewayEntryPoint}
        \nodepart{second}
          + reset(scen, seed)\\
          + step(hostIdx)\\
          + selectBaselineAction()};

      % ── Hàng 2: SimulationManager (Trung tâm) & Các VmAllocPolicy ──
      \node[cls, below=14mm of gw] (mgr)
        {\textbf{SimulationManager}
        \nodepart{second}
          + resetSimulation()\\
          + step(hostIdx)\\
          + getActionMask()\\
          + exportMetrics(dir)};
      
      \node[cls, left=8mm of mgr] (k8s)
        {\textbf{VmAllocPolicy}\\\textbf{K8sDefault}
        \nodepart{second}
          + selectHostForTask()\\
          + leastRequestedScore()};
          
      \node[cls, right=8mm of mgr] (rnd)
        {\textbf{VmAllocPolicy}\\\textbf{Random}
        \nodepart{second}
          + selectHostForTask()};

      % ── Hàng 3: Dependencies của mgr (chia làm 2 cột dưới mgr) ──
      \node[cls, below=12mm of mgr, xshift=-24mm] (fac)
        {\textbf{DatacenterFactory}
        \nodepart{second}
          + create(sim, spec, reg)\\
          + GpuState: \{total, used\}\\
          + gpuPowerWatt(state, spec)};
          
      \node[cls, below=12mm of mgr, xshift=24mm] (rdr)
        {\textbf{AlibabaTraceReader}
        \nodepart{second}
          + read(csv): List<TaskRecord>\\
          + TaskRecord (record)};

      % ── Hàng 4: Các dependencies tiếp theo đặt dưới hàng 3 ──
      % Đặt flt cách fac 16mm để có chỗ cho đường xương sống (bus) đi ngang
      \node[cls, below=16mm of fac] (flt)
        {\textbf{ScenarioFilter}
        \nodepart{second}
          + filter(tasks, scen, seed)\\
          + Scenario \{LOW,HIGH,BURST\}};
          
      % Ép MetricsExporter thẳng hàng dọc với rdr và thẳng hàng ngang với flt
      \node[cls] (exp) at (rdr |- flt)
        {\textbf{MetricsExporter}
        \nodepart{second}
          + writeCsv(path, snaps)\\
          + writeJson(path, summary)};

      % ── Config đặt bên trái của DatacenterFactory ──
      \node[cls, left=10mm of fac] (cfg)
        {\textbf{SimulationConfig}
        \nodepart{second}
          + DEFAULT\_HOST: HostSpec\\
          + DEFAULT\_POWER: PowerSpec\\
          + qosToLambda(qos): double};

      % ── Mũi tên điều hướng từ Gateway xuống Hàng 2 ──
      \coordinate (gwBus) at ($(gw.south) + (0,-6mm)$);
      \draw[arrow, gray!70] (gw.south) -- (gwBus);
      \draw[arrow] (gwBus) -| (k8s.north)
                    node[pos=0.08, above=0.5mm, font=\tiny]{uses};
      \draw[arrow] (gwBus) -- (mgr.north)
                    node[midway, right=0.5mm, font=\tiny]{wraps};
      \draw[arrow] (gwBus) -| (rnd.north)
                    node[pos=0.08, above=0.5mm, font=\tiny]{uses};

      % ── Trục xương sống từ SimulationManager xuống Hàng 3 & Hàng 4 ──
      \coordinate (mgrBus1) at ($(mgr.south) + (0,-6mm)$);
      % Lấy mốc từ đáy của DatacenterFactory (vì hộp này dài nhất) để mũi tên không đè vào hộp khác
      \coordinate (mgrBus2) at ($(fac.south) + (0,-8mm)$); 
      
      \draw[arrow, gray!70] (mgr.south) -- (mgrBus1);
      \draw[arrow] (mgrBus1) -| (fac.north);
      \draw[arrow] (mgrBus1) -| (rdr.north);
      
      % Kéo dài trục xương sống xuống hàng 4
      \coordinate (mgrBus3) at (mgrBus1 |- mgrBus2);
      \draw[arrow, gray!70] (mgrBus1) -- (mgrBus3);
      
      \draw[arrow] (mgrBus3) -| (flt.north)
                    node[pos=0.15, above=0.5mm, font=\tiny]{uses};
      \draw[arrow] (mgrBus3) -| (exp.north);

      % ── Mũi tên từ Factory sang Config ──
      \draw[arrow] (fac.west) -- (cfg.east)
                    node[midway, above=0.5mm, font=\tiny]{reads};
  \end{tikzpicture}
  \caption{Class diagram CloudSim Plus. Mũi tên đậm biểu diễn quan hệ
  \emph{depends-on}; \texttt{GatewayEntryPoint} là điểm vào duy nhất của
  Python.}
  \label{fig:class-diagram}
\end{figure}

\subsection{\texttt{SimulationConfig} --- cấu hình tập trung}
\label{ssec:config}
Mọi tham số của cluster được gom vào một file Java duy nhất, giúp khâu
tinh chỉnh và kiểm tra trở nên đơn giản. Cấu hình mặc định mô phỏng
một cụm gồm $10$ host; mỗi host có $64$ vCPU, $\SI{256}{GB}$ RAM và $8$
GPU loại A100/V100. Công suất CPU lấy ở mức $400$\,W khi đầy tải và
$120$\,W khi rảnh; GPU tương ứng là $300$\,W và $30$\,W (số liệu lấy
từ tài liệu kỹ thuật của nhà sản xuất). Bảng ánh xạ QoS $\to$ hệ số
$\lambda$ dùng trong reward SLA cũng được khai báo ở đây: bốn lớp QoS
của Alibaba (LS, Guaranteed, Burstable, BE) được ánh xạ lần lượt sang
các hệ số phạt $3{,}0$ --- $2{,}0$ --- $1{,}0$ --- $0{,}5$.

\subsection{\texttt{DatacenterFactory} và lớp theo dõi GPU}
\label{ssec:factory}
CloudSim Plus 8.5.7 chưa có khái niệm GPU. Để không phải sửa thư viện
gốc --- điều sẽ phá vỡ khả năng tái lập kết quả --- ELDAS giữ một bảng
riêng, ánh xạ mỗi host sang trạng thái GPU của nó (tổng số GPU và số
GPU đang sử dụng). Công suất GPU của host được tính theo cùng dạng
tuyến tính như CPU: số GPU đang dùng càng nhiều thì công suất càng
tiến gần mức tối đa, còn các GPU rảnh đóng góp ở mức công suất idle.
Phép tính này được gọi ở mỗi bước mô phỏng và được tính gộp vào tổng
năng lượng của cluster.

\subsection{\texttt{AlibabaTraceReader} --- đọc trace}
\label{ssec:trace}
File trace \texttt{openb\_pod\_list\_default.csv} có $11$ trường, được
ánh xạ sang một record Java đơn giản (Bảng~\ref{tab:trace-fields}). Hai
trường dẫn xuất được sinh ngay trong khi parse: \textbf{deadline} bằng
thời điểm K8s bind cộng thêm $110\%$ thời gian chạy thực tế (tức cho
phép trễ tối đa $10\%$ trước khi tính là vi phạm); \textbf{slaLambda}
là hệ số phạt theo QoS, lấy từ bảng đã khai báo ở §\ref{ssec:config}.
Những bản ghi có \texttt{podPhase = Pending} hoặc thiếu trường bắt
buộc sẽ bị bỏ qua kèm cảnh báo, không làm dừng toàn bộ quá trình
import.

\begin{table}[H]
  \centering
  \caption{Ánh xạ trường Alibaba trace $\to$ đối tượng \texttt{TaskRecord}.}
  \label{tab:trace-fields}
  \small
  \begin{tabular}{lll}
    \toprule
    \textbf{Trường CSV}      & \textbf{Trường \texttt{TaskRecord}} & \textbf{Ngữ nghĩa} \\
    \midrule
    \texttt{name}            & \texttt{name}        & Mã định danh tác vụ \\
    \texttt{cpu\_milli}      & \texttt{pesNeeded()} & $\lceil x/1000 \rceil$ PE \\
    \texttt{memory\_mib}     & \texttt{memoryMib}   & RAM yêu cầu (MiB) \\
    \texttt{num\_gpu}        & \texttt{numGpu}      & Số GPU nguyên \\
    \texttt{gpu\_milli}      & \texttt{gpuMilli}    & Tỉ phần GPU (1000 = full card) \\
    \texttt{qos}             & \texttt{slaLambda}   & $\to \lambda \in \{0{,}5,\,1,\,2,\,3\}$ \\
    \texttt{creation\_time}  & \texttt{creationTime}& Thời điểm xuất hiện (giây) \\
    \texttt{deletion\_time}  & \texttt{deletionTime}& Thời điểm kết thúc thực tế \\
    \texttt{scheduled\_time} & \texttt{scheduledTime}& Thời điểm K8s bind ban đầu \\
    \multicolumn{2}{l}{(dẫn xuất) \texttt{deadline}}& $= \mathrm{scheduledTime} + 1{,}1\,\Delta_{\mathrm{run}}$ \\
    \bottomrule
  \end{tabular}
\end{table}

\subsection{\texttt{ScenarioFilter} --- ba kịch bản tải}
\label{ssec:scenario}
Để bộc lộ ba chế độ vận hành khác biệt, \texttt{ScenarioFilter} cắt
trace đầy đủ thành ba tập con:
\begin{itemize}
  \item \textsc{Low}: $25\%$ tác vụ đầu tiên theo
        \texttt{creationTime} --- dùng để kiểm tra hành vi của
        scheduler khi cluster ở mức tải thấp (đây cũng là vùng năng
        lượng kém hiệu quả nhất theo SPECpower).
  \item \textsc{High}: toàn bộ tác vụ có
        \texttt{podPhase $\neq$ Pending} --- kịch bản tham chiếu,
        đủ khả năng làm bão hoà cluster.
  \item \textsc{Burst}: chỉ giữ những tác vụ thuộc các cửa sổ
        $1$\,giờ có cường độ đến vượt percentile $80$ --- dùng để
        kiểm tra hành vi của cluster khi bị \emph{shock} đột ngột.
\end{itemize}
Cả ba kịch bản đều dùng chung một \texttt{seed} cho \texttt{Random},
qua đó đảm bảo các lần chạy đều cho kết quả lặp lại được.

\subsection{\texttt{SimulationManager} --- ghép mô phỏng vào vòng lặp RL}
\label{ssec:sim-manager}

Đây là lớp phức tạp nhất ở Lớp 2: lớp này có nhiệm vụ điều phối nhịp
chạy của mô phỏng với nhịp ra quyết định của agent.
Hình~\ref{fig:lifecycle} mô tả luồng trạng thái trong một episode. Hai
thread --- một chạy mô phỏng phía Java, một nhận lệnh từ Python ---
trao đổi với nhau qua hai hàng đợi đồng bộ: một hàng để Python gửi
hành động sang, một hàng để Java trả lại quan sát và phần thưởng. Hai
hàng đợi được thiết kế sao cho mỗi lần \texttt{step()} là một giao
dịch trọn vẹn --- không cần khoá tường minh mà vẫn loại bỏ được hiện
tượng đua giữa hai thread.

\begin{figure}[H]
  \centering
  \begin{tikzpicture}[
      font=\scriptsize, node distance=8mm and 12mm,
      state/.style={draw, rounded corners=2pt, minimum width=22mm,
                    minimum height=8mm, align=center, fill=blue!8},
      term/.style ={draw, ellipse, minimum width=18mm, fill=green!10},
      arrow/.style={-{Latex[length=2mm]}, thick}
    ]
    \node[term] (init) {Initial};
    \node[state, right=of init]  (build) {buildSim\\(DC, broker, trace)};
    \node[state, right=of build] (emit)  {emit obs\_0\\(reward = 0)};
    \node[state, below=of emit]  (wait)  {actionQueue\\.take()};
    \node[state, left=of wait]   (apply) {allocateTask +\\advanceEnergy};
    \node[state, left=of apply]  (rew)   {compute\\Reward};
    \node[state, below=of rew]   (snap)  {recordSnapshot\\++currentTaskIdx};
    \node[state, right=of snap]  (push)  {stepResultQueue\\.put()};
    \node[term,  right=of push]  (done)  {Episode done};

    \draw[arrow] (init)--(build);
    \draw[arrow] (build)--(emit);
    \draw[arrow] (emit)--(wait);
    \draw[arrow] (wait)--(apply);
    \draw[arrow] (apply)--(rew);
    \draw[arrow] (rew)--(snap);
    \draw[arrow] (snap)--(push);
    \draw[arrow] (push) -| node[below,pos=0.25,font=\tiny]
                          {if more tasks} (wait);
    \draw[arrow] (push)--(done);
  \end{tikzpicture}
  \caption{Vòng đời \texttt{SimulationManager} trong một episode. Hai vùng
  ``actionQueue.take'' và ``stepResultQueue.put'' là hai điểm đồng bộ giữa
  thread mô phỏng và thread Py4J.}
  \label{fig:lifecycle}
\end{figure}

\textbf{Một lưu ý thiết kế quan trọng.} Trong Phase 1, ELDAS
\textbf{không khởi động đồng hồ của CloudSim}. Thay vào đó,
\texttt{SimulationManager} tự theo dõi tài nguyên đã cấp (CPU và RAM)
qua hai bảng băm riêng. Lý do: nếu bật đồng hồ mô phỏng khi chưa đăng
ký đầy đủ các listener, JVM có thể crash âm thầm do sự kiện không có
handler tương ứng. Hệ quả là Phase 1 chưa đo được thời gian hoàn thành
theo đồng hồ mô phỏng --- mọi deadline đều được suy ra từ trace ---
hạn chế này được nêu rõ ở §5.5 và sẽ được xử lý ngay đầu Phase 2.

\textbf{Reward được tính ngay tại phía Java}, theo đúng mô hình đã
phát biểu ở Chương~\ref{chap:gioi-thieu}. Phần năng lượng được tính
bằng công suất CPU và GPU hiện tại của host nhân với độ dài một bước
scheduling (đảo dấu để chuyển thành reward). Phần SLA được tính từ
mức độ trễ ước lượng vượt deadline, nhân với hệ số $\lambda$ theo
QoS. Cả hai giá trị này được trả về dưới dạng mảng hai phần tử để giữ
nguyên tính chất \textit{vector reward} của MOMDP.

% ============================================================
\section{Thiết kế Lớp 3 — Py4J Gateway}
\label{sec:layer-py4j}

\subsection{Bề mặt API tối giản}
Lớp Gateway \emph{cố ý} chỉ phơi ra một bề mặt rất nhỏ: chỉ $9$ phương
thức công khai và không trả về trực tiếp bất kỳ đối tượng CloudSim
nào. Bảng~\ref{tab:gw-api} liệt kê toàn bộ contract.

\begin{table}[H]
  \centering
  \caption{API công khai của \texttt{GatewayEntryPoint}. Mọi tham số/giá
  trị trả về đều là kiểu nguyên thuỷ, mảng nguyên thuỷ, hoặc Java
  \texttt{record} đơn giản -- thân thiện với Py4J và serializable cô lập.}
  \label{tab:gw-api}
  \footnotesize
  \begin{tabularx}{\linewidth}{l l X}
    \toprule
    \textbf{Phương thức} & \textbf{Trả về} & \textbf{Mô tả} \\
    \midrule
    \texttt{reset(scenario, seed)} & \texttt{StepResult} & Khởi tạo lại
        simulation, trả về quan sát ban đầu.\\
    \texttt{step(hostIndex)}        & \texttt{StepResult} & Áp hành động,
        trả về $(s_{t+1}, \mathbf{r}_t, d_t)$.\\
    \texttt{getActionMask()}        & \texttt{boolean[]}  & Mask hợp lệ cho
        bước hiện tại.\\
    \texttt{getActionSize()}        & \texttt{int}        & $|\mathcal{A}|$
        = số host.\\
    \texttt{getObservationSize()}   & \texttt{int}        & $6H + 4$.\\
    \texttt{getTotalEnergyKwh()}    & \texttt{double}     & Năng lượng tích
        luỹ (kWh) -- dùng cho early-stop.\\
    \texttt{selectBaselineAction(p)}& \texttt{int}        & Chỉ số host do
        baseline (\texttt{firstfit/bestfit/roundrobin/k8s/random}) chọn.\\
    \texttt{exportMetrics(dir)}     & \texttt{void}       & Ghi
        \texttt{metrics.csv} + \texttt{summary.json}.\\
    \texttt{shutdown()}             & \texttt{void}       & Dừng simulation,
        không tắt \texttt{GatewayServer}.\\
    \bottomrule
  \end{tabularx}
\end{table}

\textbf{Quy ước bất biến của hợp đồng:} mọi state đều được giữ trong
\texttt{SimulationManager} phía Java; phía Python \emph{không cache}
trạng thái. Mỗi \texttt{StepResult} là một \texttt{record} bất biến;
sau khi nhận, Python chỉ thực hiện đúng một lần chuyển đổi sang
\texttt{numpy.ndarray} (kèm thao tác clip vào $[0,1]$ để chống nhiễu
số học) rồi đẩy thẳng cho agent. Cách làm này loại bỏ triệt để lớp
bug đồng bộ trạng thái Java/Python.

\subsection{Khởi động Gateway và kết nối client}
Phía Java, \texttt{GatewayServer} được bind vào địa chỉ
\texttt{0.0.0.0} (không phải \texttt{localhost}) để container Python
trong cùng mạng Docker có thể kết nối tới được. Cổng mặc định được
đọc từ biến môi trường \texttt{PY4J\_PORT} (giá trị $25333$).

Phía Python, môi trường mở kết nối tới gateway ngay khi khởi tạo và
có cơ chế \textbf{retry với thời gian chờ tăng dần}, nhằm xử lý
khoảng vài giây ngắt quãng giữa thời điểm Docker đánh dấu Java đã
healthy và thời điểm cổng thực sự sẵn sàng. Sau $10$ lần thử mà vẫn
không kết nối được, hệ thống báo lỗi rõ ràng cho người dùng. Tham số
\texttt{auto\_convert=True} của Py4J giúp tự ép kiểu giữa Python và
Java: \texttt{int} truyền thẳng, mảng \texttt{double[]} của Java về
phía Python tự thành \texttt{list} --- nhờ đó lượng boilerplate ở phía
agent giảm đáng kể.

% ============================================================
\section{Thiết kế Lớp 4 — MO-Gymnasium Environment}
\label{sec:layer-mogym}

\subsection{Observation space}
\label{ssec:obs-space}
Quan sát là một vector phẳng dài $6H + 4$, bao gồm: $H$ giá trị CPU
util, $H$ giá trị RAM util, $H$ giá trị GPU util, $3H$ phần tử one-hot
biểu diễn trạng thái host (\textsc{Suspended}, \textsc{Idle},
\textsc{Active}) và $4$ task feature. Giá trị đã được Java chuẩn hoá
về $[0,1]$ và Python clip lại để chống lỗi số học
(Hình~\ref{fig:obs}). Phần one-hot trạng thái host là chỉ báo không
thể thiếu: nếu bỏ đi, agent không thể phân biệt host \textsc{Idle}
(sẵn sàng nhận task miễn phí) với host \textsc{Suspended} (phải trả
thêm chi phí wake-up khi đặt task lên).

\begin{figure}[H]
  \centering
  \begin{tikzpicture}[
      font=\scriptsize, node distance=0mm,
      cell/.style={draw, minimum width=11mm, minimum height=7mm, align=center},
      sub/.style ={draw, minimum width=3.4mm, minimum height=7mm, align=center,
                   inner sep=0.4mm, font=\tiny},
      cpu/.style ={cell, fill=blue!10},
      mem/.style ={cell, fill=green!10},
      gpu/.style ={cell, fill=orange!10},
      stsus/.style={sub,  fill=gray!30},
      stidl/.style={sub,  fill=yellow!25},
      stact/.style={sub,  fill=violet!20},
      task/.style={cell, fill=red!12}
    ]
    % CPU block
    \node[cpu] (c0) {$u^{\mathrm{cpu}}_1$};
    \node[cpu, right=of c0] (c1) {$u^{\mathrm{cpu}}_2$};
    \node[cpu, right=of c1] (cd) {\dots};
    \node[cpu, right=of cd] (cH) {$u^{\mathrm{cpu}}_H$};
    % MEM
    \node[mem, right=2mm of cH] (m0) {$u^{\mathrm{ram}}_1$};
    \node[mem, right=of m0] (md) {\dots};
    \node[mem, right=of md] (mH) {$u^{\mathrm{ram}}_H$};
    % GPU
    \node[gpu, right=2mm of mH] (g0) {$u^{\mathrm{gpu}}_1$};
    \node[gpu, right=of g0] (gd) {\dots};
    \node[gpu, right=of gd] (gH) {$u^{\mathrm{gpu}}_H$};

    % STATE one-hot — host 1
    \node[stsus, right=2mm of gH] (s1a) {S};
    \node[stidl, right=of s1a]    (s1b) {I};
    \node[stact, right=of s1b]    (s1c) {A};
    \node[sub,   right=of s1c, draw=none] (sd) {\dots};
    % STATE one-hot — host H
    \node[stsus, right=of sd] (sHa) {S};
    \node[stidl, right=of sHa] (sHb) {I};
    \node[stact, right=of sHb] (sHc) {A};

    % TASK (next row)
    \node[task, below=10mm of c0]      (t1) {$\hat r^{\mathrm{cpu}}$};
    \node[task, right=of t1]           (t2) {$\hat r^{\mathrm{ram}}$};
    \node[task, right=of t2]           (t3) {$\hat r^{\mathrm{gpu}}$};
    \node[task, right=of t3]           (t4) {$\lambda/3$};

    % braces
    \draw[decorate, decoration={brace,amplitude=4pt}, thick]
      (c0.north west) -- node[above=4pt,font=\scriptsize]{$H$ host CPU} (cH.north east);
    \draw[decorate, decoration={brace,amplitude=4pt}, thick]
      (m0.north west) -- node[above=4pt,font=\scriptsize]{$H$ RAM}      (mH.north east);
    \draw[decorate, decoration={brace,amplitude=4pt}, thick]
      (g0.north west) -- node[above=4pt,font=\scriptsize]{$H$ GPU}      (gH.north east);
    \draw[decorate, decoration={brace,amplitude=4pt}, thick]
      (s1a.north west) -- node[above=4pt,font=\scriptsize]{$3H$ state one-hot (\textsc{S}/\textsc{I}/\textsc{A})}
      (sHc.north east);
    \draw[decorate, decoration={brace,amplitude=4pt}, thick]
      (t4.south east) -- node[below=4pt,font=\scriptsize]{4 task features} (t1.south west);
  \end{tikzpicture}
  \caption{Cấu trúc rút gọn vector quan sát $(6H + 4)$ chiều: ba khối
  $H$ chiều cho CPU/RAM/GPU utilisation, một khối $3H$ chiều cho host
  state one-hot (\textbf{S}USPENDED, \textbf{I}DLE, \textbf{A}CTIVE
  cho mỗi host), và bốn ô cuối mô tả tác vụ kế tiếp cần xếp lịch.}
  \label{fig:obs}
\end{figure}

Việc \emph{nhóm theo loại tài nguyên} (CPU, rồi RAM, rồi GPU, rồi
state) thay vì \emph{nhóm theo host} là một lựa chọn có chủ đích.
Lớp ẩn của MLP học các tương tác \emph{trong cùng một loại} dễ hơn
(ví dụ: host nào có CPU thấp nhất, host nào đang \textsc{Suspended}).
Sau này, nếu chuyển sang GNN ở stretch goal, các khối CPU/RAM/GPU có
thể \texttt{reshape} về ma trận $H \times 3$ một cách tự nhiên; còn
khối state được đưa vào node feature dưới dạng vector one-hot $3$
chiều.

\subsection{Action space và action masking}
\label{ssec:action-space}
Không gian hành động khá đơn giản: là chỉ số của host (từ $0$ đến
$H - 1$) mà agent muốn đặt tác vụ hiện tại lên. Ở mỗi bước, môi
trường gọi sang Java để lấy về một mask boolean dài $H$: phần tử thứ
$j$ bằng \texttt{true} nếu host $j$ còn đủ CPU, RAM và GPU cho tác
vụ; ngược lại là \texttt{false}. MaskablePPO áp mask này lên đầu ra
của policy, qua đó đảm bảo agent không bao giờ chọn host không đủ
tài nguyên. Trong trường hợp đặc biệt khi không có host nào chứa nổi
tác vụ (toàn bộ mask đều \texttt{false}), phía Python sẽ bật lại
toàn bộ mask để tránh agent bị ``khoá cứng''.

\subsection{Vector reward và linear scalarization}
\label{ssec:reward}
Ở mỗi bước, môi trường trả về reward dưới dạng vector hai phần tử:
phần năng lượng (luôn âm, bằng lượng năng lượng tiêu thụ thêm trong
bước đó, đã đảo dấu) và phần SLA (âm nếu vi phạm deadline, $0$ nếu
đúng hạn). Để có thể chạy được với PPO chuẩn, ELDAS cung cấp một
wrapper là \texttt{ScalarRewardWrapper}, có nhiệm vụ cộng hai thành
phần này theo trọng số $(w_{\mathrm{e}},\,w_{\mathrm{s}})$ với
$w_{\mathrm{e}} + w_{\mathrm{s}} = 1$. Trọng số mặc định là
$(0{,}8;\,0{,}2)$; Phase 2 sẽ quét $w_{\mathrm{e}} \in \{0{,}1;\,
\dots;\, 0{,}9\}$ để dựng Pareto front. Khi cần benchmark với một
thuật toán MORL chuyên dụng, chỉ cần bỏ wrapper là môi trường trả
thẳng vector reward về cho thuật toán --- không phải đụng đến code
Java.

\subsection{Điều kiện kết thúc episode}
\texttt{terminated\,=\,True} chỉ khi mọi tác vụ trong tập
\texttt{tasks} đã được phân bổ
(\texttt{currentTaskIdx == tasks.size()}). \texttt{truncated} luôn
bằng \texttt{False} ở Phase 1, do bài toán có chân trời hữu hạn một
cách tự nhiên (số tác vụ cố định sau khi \texttt{ScenarioFilter}
cắt). Sang Phase 2, ELDAS có thể bật \texttt{truncated\,=\,True} khi
số bước vượt \texttt{max\_steps}, qua đó giới hạn chi phí cho các
chính sách tồi.

\subsection{Tóm tắt vòng lặp}
Một episode hoàn chỉnh theo đúng chuẩn Gymnasium gồm bốn bước lặp lại:
(1) gọi \texttt{action\_masks()} để biết các host khả thi; (2) policy
chọn một host hợp lệ; (3) gọi \texttt{step(action)} để nhận quan sát
và reward mới; (4) kiểm tra điều kiện kết thúc. Sau khi episode kết
thúc, hàm \texttt{export\_metrics()} được gọi để Java ghi
\texttt{metrics.csv} và \texttt{summary.json} ra volume kết quả.

% ============================================================
\section{Thiết kế Lớp 1 — Docker Compose Infrastructure}
\label{sec:layer-docker}

\subsection{Bố trí container và volume}
Hai container được nối qua mạng nội bộ \texttt{sim-net} và chia sẻ ba
volume (Hình~\ref{fig:docker-topo}). Volume chứa trace được mount ở
chế độ \textbf{chỉ đọc} cho \texttt{cloudsim-java}, qua đó Java không
có quyền ghi đè dữ liệu thực nghiệm và tránh được tình huống vô tình
làm hỏng trace gốc trong khi debug.

\begin{figure}[H]
  \centering
  \begin{tikzpicture}[
      font=\scriptsize, node distance=8mm and 14mm,
      svc/.style ={draw, rounded corners=2pt, minimum width=42mm,
                   minimum height=18mm, align=center, fill=blue!8},
      vol/.style ={draw, rounded corners=2pt,
                   minimum width=22mm, minimum height=10mm,
                   align=center, fill=yellow!25},
      net/.style ={draw, dashed, rounded corners=2pt,
                   inner sep=4mm, fill=gray!5},
      arrow/.style={-{Latex[length=2mm]}, thick}
    ]
    \node[svc] (java) {\textbf{cloudsim-java}\\(JRE 21, $\sim$380\,MB)\\
                       \texttt{:25333}};
    \node[svc, right=22mm of java] (py) {\textbf{rl-agent}\\
                       (Python 3.11.9)\\client only};
    \begin{pgfonlayer}{background}
      \node[net, fit=(java)(py),
            label={[anchor=north,yshift=-1mm,font=\scriptsize\itshape]
                   north:bridge \texttt{sim-net}}] (net) {};
    \end{pgfonlayer}

    \node[vol, below=of java]               (vt) {trace-data\\(:ro)};
    \node[vol, below=of $(java.south)!0.5!(py.south)$] (vr) {results};
    \node[vol, below=of py]                 (vm) {model-store};

    \draw[arrow] (java) -- (vt);
    \draw[arrow] (java) -- (vr);
    \draw[arrow] (py)   -- (vr);
    \draw[arrow] (py)   -- (vm);
    \draw[arrow, {Latex[length=2mm]}-{Latex[length=2mm]}]
      (py.west) to[bend right=10] node[above,font=\tiny]
        {Py4J TCP :25333} (java.east);
  \end{tikzpicture}
  \caption{Topology Docker Compose. \texttt{rl-agent} không expose port
  nào (client thuần); \texttt{cloudsim-java} chỉ expose
  \texttt{\$PY4J\_PORT}.}
  \label{fig:docker-topo}
\end{figure}

\subsection{Multi-stage build và tối ưu kích thước image}
Cả hai Dockerfile đều dùng pattern \textbf{multi-stage}: một stage
chứa các công cụ biên dịch nặng (Maven, gcc, \dots), stage còn lại
chỉ giữ phần runtime tối thiểu. Nhờ vậy, image
\texttt{cloudsim-java} giảm từ khoảng $1{,}1$\,GB xuống chỉ còn
$\sim 380$\,MB; image \texttt{rl-agent} cũng giảm tương tự so với
base \texttt{python:3.11}. Có một thủ thuật nhỏ nhưng đáng giá: copy
các file khai báo dependency (\texttt{pom.xml},
\texttt{requirements.txt}) \emph{trước} khi copy source code --- khi
các file này không đổi, Docker sẽ tái sử dụng layer cache và mỗi lần
build lại tiết kiệm được $1$--$2$ phút.

\subsection{Healthcheck và thứ tự khởi động}
Để tránh hiện tượng đua khi khởi động, \texttt{cloudsim-java} có một
healthcheck kiểm tra cổng Py4J ($25333$) đã lắng nghe hay chưa;
song song, \texttt{rl-agent} được khai báo phải chờ Java sẵn sàng
mới được chạy. Docker Compose đảm nhiệm phần điều phối này. Cơ chế
retry ở phía Python (đã mô tả ở §\ref{sec:layer-py4j}) đóng vai trò
lớp phòng vệ thứ hai, xử lý khoảng vài giây ngắn giữa thời điểm cổng
mở và thời điểm gateway hoàn tất bind.

\subsection{Khả năng tái lập kết quả}
ELDAS được thiết kế để bảo đảm \textbf{tái lập kết quả}: với cùng
\texttt{RANDOM\_SEED}, mỗi lần chạy đều cho cùng cách chia tác vụ
giữa các kịch bản, cùng thứ tự lấy mẫu của các baseline ngẫu nhiên
và cùng một chuỗi hành động. Mọi phụ thuộc đều được ghim phiên bản
cứng: phía Java có Java 21 (eclipse-temurin), Maven 3.9.9, CloudSim
Plus 8.5.7 và Py4J 0.10.9.7; phía Python có Python 3.11.9,
mo-gymnasium 1.3.2 và stable-baselines3 2.7.0. Riêng
\texttt{gymnasium} cố ý không pin, để pip tự giải xung đột giữa SB3
và MO-Gym.

% ============================================================
\section{Thiết kế Lớp 5 — Observability stack (tùy chọn)}
\label{sec:layer-observability}

Bốn lớp ở trên đã đủ để chạy huấn luyện và xuất số liệu offline. Tuy
vậy, khi cần \emph{quan sát chiến lược đặt task của scheduler ngay
trong lúc chạy} --- chẳng hạn để minh hoạ tương phản giữa ``K8s rải
đều'' và ``MORL gom chặt'' trong báo cáo hay buổi bảo vệ --- thì các
biểu đồ tĩnh từ \texttt{metrics.csv} không đủ trực quan. ELDAS vì thế
bổ sung thêm một \textbf{lớp Observability tuỳ chọn} (Phase 1.5) gồm
Prometheus + Grafana, được thiết kế dưới bốn ràng buộc:

\begin{enumerate}
  \item \textbf{Không can thiệp vào đường tới hạn (hot path).} Vòng
        lặp RL phải chạy bình thường ngay cả khi lớp này bị tắt hoặc
        gặp lỗi. Bởi vậy, Prometheus và Grafana \emph{không} được
        đưa vào \texttt{depends\_on} của bất kỳ service nào.
  \item \textbf{Bật/tắt qua Docker Compose profile.} Mặc định,
        \texttt{docker compose up} chỉ khởi động hai container core;
        chỉ khi chạy \texttt{docker compose -{}-profile monitoring up}
        thì stack quan sát mới được bật lên.
  \item \textbf{Java là nguồn metric chính.}
        \texttt{cloudsim-java} là service \emph{long-running} có
        healthcheck, nên Prometheus có thể scrape ổn định; ngược
        lại, \texttt{rl-agent} là \emph{ephemeral} (chạy bằng
        \texttt{docker compose run -{}-rm}) --- nếu phát metric từ
        Python thì target sẽ chập chờn. Mặt khác, phía Java giữ các
        giá trị \emph{thực} (kWh, queue length), còn Python chỉ
        thấy observation đã được chuẩn hoá về $[0,1]$.
  \item \textbf{Không thay thế kết quả khoa học.} Grafana scrape
        theo \emph{wall-time} chứ không phải \emph{sim-time}, do đó
        trục thời gian của Grafana không tương thích với báo cáo
        định lượng. Vì lý do này, các figure trong
        Chương~\ref{chap:ket-qua} vẫn lấy từ \texttt{metrics.csv}
        và WandB; Grafana chỉ phục vụ mục đích debug và demo trực
        tiếp.
\end{enumerate}

\subsection{Sơ đồ tích hợp}
Hình~\ref{fig:obs-topo} minh hoạ cách hai service mới gắn vào mạng
\texttt{sim-net} mà vẫn không chạm vào đường dữ liệu Py4J giữa hai
container core.

\begin{figure}[H]
  \centering
  \begin{tikzpicture}[
      font=\scriptsize, node distance=8mm and 10mm,
      svc/.style ={draw, rounded corners=2pt, minimum width=34mm,
                   minimum height=16mm, align=center, fill=blue!8},
      obs/.style ={draw, rounded corners=2pt, minimum width=34mm,
                   minimum height=16mm, align=center, fill=green!10},
      net/.style ={draw, dashed, rounded corners=2pt,
                   inner sep=4mm, fill=gray!5},
      arrow/.style={-{Latex[length=2mm]}, thick},
      scrape/.style={-{Latex[length=2mm]}, dashed, thick, gray!70}
    ]
    \node[svc] (java) {\textbf{cloudsim-java}\\
                       \texttt{:25333} (Py4J)\\
                       \texttt{:9091} (\textit{exporter})};
    \node[svc, right=22mm of java] (py) {\textbf{rl-agent}\\
                       (ephemeral)\\
                       \texttt{:8000} (P2)};
    \node[obs, below=14mm of java] (prom)
                     {\textbf{prometheus}\\
                      \texttt{:9090}\\
                      scrape 2\,s};
    \node[obs, below=14mm of py] (graf)
                     {\textbf{grafana}\\
                      \texttt{:3000}\\
                      auto-provision};
    \begin{pgfonlayer}{background}
      \node[net, fit=(java)(py)(prom)(graf),
            label={[anchor=north,yshift=-1mm,font=\scriptsize\itshape]
                   north:bridge \texttt{sim-net}}] (net) {};
    \end{pgfonlayer}
    \draw[arrow, {Latex[length=2mm]}-{Latex[length=2mm]}]
      (py.west) to[bend right=8] node[above,font=\tiny]
        {Py4J} (java.east);
    \draw[scrape] (prom.north) -- node[right,font=\tiny] {/metrics} (java.south);
    \draw[scrape] (prom.east) -- node[above,font=\tiny] {/metrics (P2)} (graf.west);
    \draw[scrape] (graf.south west) to[bend right=10] node[below,font=\tiny] {PromQL} (prom.east);
  \end{tikzpicture}
  \caption{Topology lớp Observability. Mũi tên liền là dữ liệu trên đường
  tới hạn (Py4J); mũi tên đứt là metric scrape, chạy ngoài đường tới hạn.
  Hai service \texttt{prometheus} và \texttt{grafana} thuộc Docker
  Compose profile \texttt{monitoring} -- mặc định không khởi động.}
  \label{fig:obs-topo}
\end{figure}

\subsection{Mô hình dữ liệu metric}
Java exporter đăng ký chín họ metric của ELDAS, chia thành ba nhóm
theo ngữ nghĩa của Prometheus (Bảng~\ref{tab:metric-model}). Tất cả
đều được đính kèm hai tag bối cảnh \texttt{scenario} và
\texttt{scheduler}, nhờ vậy trên cùng một dashboard có thể so sánh
trực tiếp \emph{cùng một tình huống tải nhưng khác chiến lược lập
lịch}.

\begin{table}[H]
  \centering
  \caption{Mô hình dữ liệu metric của ELDAS Java exporter.}
  \label{tab:metric-model}
  \footnotesize
  \begin{tabularx}{\linewidth}{l l X}
    \toprule
    \textbf{Họ metric} & \textbf{Kiểu} & \textbf{Ngữ nghĩa} \\
    \midrule
    \texttt{eldas\_host\_\{cpu,mem,gpu\}\_util}
        & Gauge & Tỉ lệ sử dụng từng host $\in[0,1]$.
                 Label: \texttt{host\_id} (0..9). \\
    \texttt{eldas\_host\_power\_watt}
        & Gauge & Công suất tức thời từng host (W). \\
    \texttt{eldas\_total\_energy\_kwh}
        & Gauge & Năng lượng tích lũy trong episode đang chạy. Reset về
                 $0$ mỗi \texttt{resetSimulation()}. \\
    \texttt{eldas\_pending\_tasks}
        & Gauge & Số task còn lại trong episode -- chỉ báo
                 backpressure. \\
    \texttt{eldas\_sim\_clock\_seconds}
        & Gauge & Đồng hồ mô phỏng. Đối chiếu wall-clock để phát hiện
                 stall. \\
    \texttt{eldas\_sla\_violations\_total}
        & Counter & Số vi phạm SLA \emph{từ lúc JVM khởi động}.
                   Không reset; truy vấn bằng \texttt{rate()} /
                   \texttt{increase()}. Label phụ:
                   \texttt{qos} $\in$ \{BE, Burstable, Guaranteed, LS\}. \\
    \texttt{eldas\_tasks\_scheduled\_total}
        & Counter & Số placement thành công, cộng dồn. \\
    \bottomrule
  \end{tabularx}
\end{table}

\paragraph{Vì sao \texttt{episode} không được dùng làm label?}
Đây là một chi tiết cần lưu ý khi thiết kế. Nếu để
\texttt{episode\_id} làm label cho gauge ở mức từng host,
cardinality của metric sẽ phình ra theo số episode trong weight
sweep ($\sim 10\text{ host}\times 3\text{ scenario}\times
4\text{ scheduler} \times N\text{ episode}$, với $N$ có thể từ
$100$ trở lên). Cách làm đúng, theo khuyến nghị của Prometheus, là
giữ cho cardinality bị chặn (\textit{bounded}): coi
\emph{Prometheus time} là chiều thời gian, còn label chỉ dùng để
\emph{định danh thực thể}. Mỗi lần \texttt{resetSimulation()}
được gọi, gauge sẽ được đưa về $0$, nhờ vậy các điểm gãy của
gauge tự đánh dấu ranh giới giữa các episode ngay trên trục thời
gian.

\subsection{Bốn nguyên tắc an toàn}
Đặc tính then chốt của lớp này là \textbf{không được gây hại}. Bốn
nguyên tắc dưới đây được hiện thực ngay ở mức mã nguồn (chi tiết tại
§\ref{sec:impl-monitoring}):

\begin{itemize}
  \item \emph{No-op khi tắt.} Cờ \texttt{enabled} mặc định bằng
        \texttt{false}; khi đó mọi setter chỉ thực hiện đúng một
        lệnh \texttt{return}, không cấp phát object và không I/O.
  \item \emph{Mọi mutator đều được bọc trong
        \texttt{try/catch Throwable}.} Lỗi bên monitoring tuyệt
        đối không thể lan ra làm crash JVM sim.
  \item \emph{Idempotent.} Hàm \texttt{MetricsRegistry.start()}
        có thể được gọi nhiều lần một cách an toàn; \texttt{stop()}
        làm sạch \texttt{CollectorRegistry} để có thể tái sử dụng
        cổng.
  \item \emph{Counter chỉ monotonic.} Đúng theo semantic của
        Prometheus --- không bao giờ \emph{reset} counter trong
        cùng một vòng đời JVM; nếu cần reset, dùng gauge là cách
        thay thế hợp lệ.
\end{itemize}

% ============================================================
\section{Tổng kết chương}
\label{sec:ket-luan-c3}

Chương này đã trình bày kiến trúc bốn lớp của ELDAS, kèm theo lý do
lựa chọn từng công nghệ cốt lõi. Bảng~\ref{tab:requirement-trace}
ánh xạ ngược các mục tiêu đã nêu ở Chương~\ref{chap:gioi-thieu} sang
thành phần kiến trúc chịu trách nhiệm tương ứng, qua đó cho thấy mọi
mục tiêu đều đã có vị trí cụ thể trong thiết kế.

\begin{table}[H]
  \centering
  \caption{Ánh xạ ngược từ mục tiêu Phase 1 (§\ref{sec:muc-tieu}) sang thành phần kiến trúc.}
  \label{tab:requirement-trace}
  \footnotesize
  % Đảm bảo cột X ở giữa để nó tự động co giãn hết phần không gian còn lại của \linewidth
  \begin{tabularx}{\linewidth}{l l X}
    \toprule
    \textbf{Mục tiêu} & \textbf{Mô tả ngắn} & \textbf{Thành phần} \\
    \midrule
    MT1 & Engine mô phỏng + năng lượng + GPU tracking
        & §\ref{sec:layer-cloudsim} (Lớp 2) \\
    MT2 & Tích hợp Alibaba trace + 3 kịch bản
        & §\ref{ssec:trace}, §\ref{ssec:scenario} \\
    MT3 & Cầu Java--Python với pause/resume/reset
        & §\ref{sec:layer-py4j} (Lớp 3) + §\ref{ssec:sim-manager} \\
    MT4 & MO-Gym env + vector reward + action masking
        & §\ref{sec:layer-mogym} (Lớp 4) \\
    MT5 & Năm baseline cổ điển + MORL PPO tối thiểu
        & \texttt{VmAllocationPolicy}\{FirstFit, BestFit, RoundRobin, K8sDefault, Random\} + \texttt{train\_min.py} \\
    MT6 & Đóng gói Docker Compose, bảo đảm tái lập
        & §\ref{sec:layer-docker} (Lớp 1) \\
    \bottomrule
  \end{tabularx}
\end{table}

Có thể tóm lược ba điểm nổi bật của kiến trúc như sau:
(i)~\textbf{tách rời rõ giữa các lớp} --- có thể thay engine mô phỏng
hoặc thuật toán RL một cách độc lập, không phải động đến phần còn lại;
(ii)~\textbf{vòng lặp RL chạy nguyên tử} --- hai hàng đợi đồng bộ phía
Java cùng cơ chế action mask phía Python loại bỏ hiện tượng đua mà
không cần đến khoá tường minh; (iii)~\textbf{khả năng tái lập kết
quả} --- nhờ seed chung, các phiên bản được ghim cứng và Docker image
bất biến, chạy lại trên máy khác vẫn ra cùng một bộ số.

Chương kế tiếp sẽ trình bày chi tiết hiện thực ở mức code: cách parse
CSV, walk-through năm baseline scheduler, hành vi của
\texttt{MetricsExporter} và kết quả smoke-test xác nhận vòng RL hoàn
chỉnh qua bốn lớp đã chạy đúng.
